%=======================================================================
% Requires: booktabs, amsmath, amssymb, siunitx, todonotes (\todo, \unsure)
% Bib keys used: s3e, cumulti, deepsense6g, inscope  (see refs.bib)
%=======================================================================
\section{Dataset Design and Collection}
\label{sec:collection}

\subsection{Environments}
\label{sec:environments}
% Indoor/outdoor sites chosen for viewpoint, appearance, and radio diversity.

We collected data at National Yang Ming Chiao Tung University (NYCU) Guangfu
Campus, across indoor and outdoor sites. Sites were selected so that the three
factors governing collaborative perception vary across the release, and each is
characterised by a quantity we measure rather than by a description:

\begin{itemize}
\item \emph{Viewpoint.} Whether an elevated, persistent vantage exists, and how
far it can be displaced from the mobile nodes \emph{in the ground plane}. What
matters is the range of aspect angles $\beta$ a site can produce between an ego
and a helper, not the mast height alone; Section~\ref{sec:results} shows the two
are not interchangeable.

\item \emph{Appearance.} Illumination, texture, and clutter density, which set
how much evidence a node's front end returns at a given range---and therefore
where the gap between \emph{visible} and \emph{observed} opens
(Eqs.~\eqref{eq:visible}--\eqref{eq:observed}).

\item \emph{Radio propagation.} Whether structure intervenes between a node and
the access point. Indoor sites obstruct the first Fresnel zone at short range;
outdoor sites hold clearance over longer links. Large-scale sensing-and-communication
datasets infer blockage from received power itself
\cite{deepsense6g}; we instead compute clearance geometrically from the map and
validate it against measured RSSI (Section~\ref{sec:results}), so link state is
an independent label rather than a threshold on the signal it is meant to
explain.
\end{itemize}

Every site is surveyed and mapped once as a dense point cloud, and all nodes are
expressed in that single anchored map frame, so occlusion, co-visibility, and
link geometry are computed against one geometry rather than against per-node
local frames. The sequence characterised in this paper was collected at the MIRC
site: an indoor laboratory with a single surveyed elevated camera--radar mast, a
\SI{5.75}{\giga\hertz} access point placed behind furniture, and static seating
that provides the occluding structure.

\todo{Site-overview figure: map cloud, surveyed infrastructure pose, AP
position, and both mobile trajectories, one panel per site.}

\subsection{Collection Protocol}

\subsubsection{Roles}
One mobile agent is designated \textit{ego} and is deliberately viewpoint
limited: its observations are degraded by design, either by trajectory (it is
routed so that structure falls between it and the annotated objects) or by
sensing budget (it carries the shorter-range depth sensor). The second mobile
agent is the \textit{partner}. Ego assignment alternates across runs so that a
reported condition is a property of the geometry rather than of a particular
platform; for the single-sequence analysis of Section~\ref{sec:results} we
report \emph{both} ego assignments of the same bag, which is the strongest
available form of this control.

Infrastructure nodes are never ego. Their role is the elevated, persistent
vantage that no mobile node can hold. Existing multi-robot and infrastructure
datasets collaborate between sensors at a common height---$\Delta h = 0$ in
S3E~\cite{s3e} and in roadside sets such as InScope~\cite{inscope}---so the
elevation asymmetry is absent by construction; treating our infrastructure node
as ego would reintroduce exactly that symmetry and collapse the asymmetry this
dataset is built to expose.

\subsubsection{Trajectory Paradigms}
\label{sec:paradigms}
A paradigm specifies a trajectory geometry \emph{and} the collaboration
condition that geometry is intended to induce. Paradigm-driven trajectory
design is established practice: S3E~\cite{s3e} names four collaborative
trajectory paradigms and grades each against loop-closure opportunity with
filled and open circles, and CU-Multi~\cite{cumulti} goes further by treating
inter-trajectory overlap as a designed, controlled variable. Neither releases a
per-frame measurement of what the design actually produced. S3E's grading is an
assertion about a run; CU-Multi's overlap is controlled by re-running a single
platform four times, so the robots are never simultaneously present and no
instantaneous collaboration geometry exists to measure. We keep the vocabulary
and replace the grading with measurement: every paradigm is defined as a
predicate over per-frame quantities that we compute from the map, the
annotations, and the poses, and that we release alongside the bags.

\paragraph{Per-observation predicates.}
Fix an annotated object and a frame. Let $e$ denote the ego node and $j$ range
over the other nodes (the peer robot and the infrastructure node). For node
$i$, let $\omega_i \in [0,1]$ be the fraction of ray casts from $i$'s sensor
origin to the object's surface that are blocked by mapped structure, $r_i$ the
range to the object, and $n_i$ the count of evidence returns falling inside the
annotated box (LiDAR points, depth points, or projected pixel area, depending
on the node's primary modality). With an occlusion threshold $\tau = 0.5$,
%
\begin{align}
\mathrm{visible}(i) &\equiv \text{in FoV} \;\wedge\; \omega_i < \tau,
\label{eq:visible}\\
\mathrm{observed}(i) &\equiv \mathrm{visible}(i) \;\wedge\; n_i \ge n_{\min,i}
\;\wedge\; r_i \le r_{\max,i}.
\label{eq:observed}
\end{align}
%
Occlusion has been quantified on infrastructure sensors before:
InScope~\cite{inscope} reports an occlusion degree $\xi_D$ over roadside LiDAR.
That measure is computed for a fixed roadside vantage, between sensors at a
common height, and is summarised per dataset rather than released per frame, so
it cannot be used to condition an evaluation. Our $\omega_i$ is defined per
(frame, node, object) and ships as a column.

The distinction between~\eqref{eq:visible} and~\eqref{eq:observed} carries most
of the analytical weight in this paper. \emph{Visible} is a geometric claim: an
unobstructed line exists. \emph{Observed} is an evidence claim: the sensor
actually returned enough of the object to support a detection, within the range
at which that modality still returns anything. A node can be visible and
unobserved for an entire sequence, and Section~\ref{sec:results} shows this is
the dominant failure mode here.

A node $j \ne e$ is a \emph{helper} for the observation if
$\mathrm{visible}(j)$; the \emph{best} helper $j^\star$ is the helper with the
lowest $\omega$. The aspect angle $\beta$ is the angle at the object between
the rays to $e$ and to $j^\star$, so $\beta$ near $0^\circ$ means the helper
adds a redundant viewpoint and $\beta > 90^\circ$ means it adds a genuinely
complementary one. Per frame we also record the co-visibility fraction $C$
(objects seen by at least two nodes, over objects seen by at least one), and,
for each link, the minimum clearance of the first Fresnel zone against the same
occupancy grid, normalised by the Fresnel radius; $\mathrm{clear}(a,b) < 0.6$
is our obstructed-link criterion.

\paragraph{Paradigm assignment.}
Table~\ref{tab:paradigms} lists the paradigms. They are evaluated in order and
the first match wins, so the ordering encodes a precedence: a condition that
combines occlusion with a degraded link is reported as such rather than being
absorbed into the plain occlusion class. $v_o$ is the object's speed, $v_{r,i}$
its radial velocity relative to node $i$, and $v_{\min}$ the Doppler resolution
floor of the radar, $v_{\min} = k\lambda / (2 N_c T_c)$ with $k \approx 1.5$.

\begin{table}[t]
\centering
\caption{Trajectory paradigms. Evaluated in order; first match wins.
$\tau=0.5$; helper and $j^\star$ as defined in Section~\ref{sec:paradigms}.
The rightmost column states what the current release can support: \checkmark{}
computed and released, $\circ$ computed but not yet folded into the classifier,
$\times$ undefined without moving targets.}
\label{tab:paradigms}
\small
\begin{tabular}{@{}clll@{}}
\toprule
\# & Paradigm & Condition & Status \\
\midrule
1 & Doppler diversity   & $\lVert v_o \rVert > 0.3 \;\wedge\; \lVert v_{r,e}\rVert < v_{\min} \;\wedge\; \max_j \lVert v_{r,j}\rVert \ge v_{\min}$ & $\times$ \\
2 & Occlusion $\times$ link & $\neg\,\mathrm{visible}(e) \;\wedge\; \text{helper} \;\wedge\; \mathrm{clear}(e,j^\star) < 0.6$ & $\circ$ \\
3 & Occlusion           & $\neg\,\mathrm{visible}(e) \;\wedge\; \text{helper}$ & \checkmark \\
4 & Handover            & $\lvert \{ i : \mathrm{visible}(i) \} \rvert = 1$ & \checkmark \\
5 & Link stress         & $\text{helper} \;\wedge\; \mathrm{clear}(e,j^\star) < 0.6$ & $\circ$ \\
6 & Co-observation      & $\mathrm{visible}(e) \;\wedge\; \text{helper}$ & \checkmark \\
7 & Unstructured        & none of the above & \checkmark \\
\bottomrule
\end{tabular}
\end{table}

\paragraph{Labels are descriptive, not prescriptive.}
Two properties of this scheme matter for how the dataset should be used.
First, a paradigm label is assigned \emph{per frame from measured metrics}, not
declared once per run: a sequence intended as co-observation contains occlusion
frames whenever the geometry produced them, and they are labelled as such.
Second, and consequently, \textbf{evaluation strata are cut by the raw metrics,
never by the label.} A method is reported against $\omega$, $\beta$, $r$, $C$,
and link clearance directly; the label exists so that a run can be described in
one word, and so that intent can be audited. Each run additionally records an
\texttt{intended\_paradigm} in its metadata, and the agreement between intended
and measured is the design check we report in Section~\ref{sec:results}: it
says how reliably a trajectory plan produces the condition it was written to
produce, which is precisely the quantity a $\bullet/\circ$ grading~\cite{s3e}
asserts without measuring.

\paragraph{Standing parameters.}
The geometric pipeline is frozen across sequences: occlusion threshold
$\tau = 0.5$; occupancy voxel \SI{10}{\centi\metre}; 200 rays per (object,
node, frame); statistical outlier removal at 20 neighbours and $2\sigma$; a
voxel counts as occupied only with $\ge 8$ points, and a ray is blocked only on
2 consecutive occupied voxels; occluder hits within \SI{0.6}{\metre} of a
sensor's own origin and occupancy within \SI{0.7}{\metre} of either link
endpoint are ignored (both nodes and the access point are themselves present in
the map); evidence floors of 1 LiDAR point, 50 depth points, and 400\,px$^2$;
analysis range caps of \SI{20}{\metre} (Ouster), \SI{4.5}{\metre} (ZED),
\SI{5}{\metre} (RealSense), and \SI{20}{\metre} (infrastructure pixels), each
calibrated empirically from the range at which the median evidence count falls
to its floor; obstructed link below $0.6\,r_1$; carrier frequency read per
frame from the WiFi status stream (\SI{5.75}{\giga\hertz}). Density
thresholding, consecutive-hit testing, and endpoint clearance are workarounds
for map noise and for fixed nodes being mapped as structure; they are recorded
here because they change the numbers and a re-analysis with a survey-grade map
should be able to drop them.

\subsection{Conditions Captured}
\label{sec:conditions}

Table~\ref{tab:inventory} inventories the current release. It is one sequence,
and we state its scale plainly rather than reporting frame counts as if they
were independent samples: \num{733} synchronised frames over \SI{156}{\second}
contain \num{9} and \num{8} occlusion episodes under the two ego assignments,
and the episode is the unit at which occlusion claims are supported.

\begin{table}[t]
\centering
\caption{Inventory of the current release. Deliberate difficulty is what the
protocol induces; incidental difficulty is what the geometry produced without
staging, and is reported separately because unstaged occurrences are the
stronger evidence.}
\label{tab:inventory}
\small
\begin{tabular}{@{}ll@{}}
\toprule
Sequences / sites & 1 / 1 (MIRC, indoor) \\
Duration, frames & \SI{156}{\second}, \num{733} frames at \SI{5}{\hertz} \\
Nodes & 2 mobile + 1 surveyed elevated infrastructure \\
Ego assignments & 2 (both mobile nodes, same bag) \\
Annotated objects & 2 static (chairs), single anchored map frame \\
Modalities & LiDAR, stereo depth, RGB-D, \SI{60}{\giga\hertz} radar (2 mobile + 1 infra), \\
          & WiFi RSSI / CSI per frame, NTP offsets \\
Released geometry & per-frame occlusion, visibility, evidence counts, $C$, $\beta$, \\
          & Fresnel clearance and link state, occlusion episodes \\
\midrule
Deliberate difficulty & viewpoint-limited ego; \SI{1.53}{\metre} elevation offset \\
          & (IQR \SIrange{1.52}{1.54}{\metre}); AP placed behind structure \\
Incidental difficulty & 9 / 8 occlusion episodes (median \SI{18.8}{\second}); \\
          & obstructed link on 33\,\% / 18\,\% of frames \\
\midrule
Not captured & scripted motion (Doppler paradigm undefined); \\
          & a second infrastructure node (handover under-sampled); \\
          & radar on \texttt{mobile\_2}; low-light and dropout conditions \\
\bottomrule
\end{tabular}
\end{table}

\todo{Extend this table once the staged runs land: scripted participants,
low-light repeat, node-outage repeat, and a bounded fully-annotated region.
Report hours, distance, frames, and annotated instances for the full release.}
